\section{数量场在曲面上的积分}

\begin{problem}{曲面面积计算}{Surface-Area-Calculations}
    求下列曲面在指定部分的面积：
    \begin{enumerate}
        \item 锥面 $z = \sqrt{x^2 + y^2}$ 包含在圆柱 $x^2 + y^2 = 2x$ 内的部分；
        \item 柱面 $x^2 + y^2 = a^2$ 被平面 $x + z = 0, x - z = 0$ ($x > 0, y > 0$) 所截的那部分；
        \item 圆柱面 $x^2 + y^2 = a^2$ 被圆柱 $y^2 + z^2 = a^2$ 所割下的那部分；
        \item 球面 $x^2 + y^2 + z^2 = 3a^2$ 和抛物面 $x^2 + y^2 = 2az$ ($z \ge 0$) 所围成的立体的全表面；
        \item 曲面 $x = \frac{1}{2}(2y^2 + z^2)$ 被柱面 $4y^2 + z^2 = 1$ 所截下的那部分；
        \item 锥面 $z^2 = x^2 + y^2$ 被 $Oxy$ 平面和 $z = \sqrt{2} \left( \frac{x}{2} + 1 \right)$ 所截下的那部分；
        \item 螺旋面 $x = r \cos \varphi, y = r \sin \varphi, z = h \varphi$ 在 $0 < r < a, 0 < \varphi < 2\uppi$ 那部分；
        \item 曲面 $(x^2 + y^2 + z^2)^2 = 2a^2 xy$ 的全部．
    \end{enumerate}
\end{problem}

\begin{solution}
    \ansitem{1}

    投影区域为 $D = \{(x, y) \mid (x-1)^2 + y^2 \le 1\}$．由 $z = \sqrt{x^2+y^2}$ 得
    \begin{equation}
        \d S = \sqrt{1 + z_x^2 + z_y^2} \d x \d y = \sqrt{1 + \frac{x^2}{x^2+y^2} + \frac{y^2}{x^2+y^2}} \d x \d y = \sqrt{2} \d x \d y
    \end{equation}
    故面积为
    \begin{equation}
        S = \iint_D \sqrt{2} \d x \d y = \sqrt{2} \cdot \operatorname{Area}(D)
    \end{equation}
    \begin{equation}
        \boxed{S = \sqrt{2} \uppi}
    \end{equation}

    \ansitem{2}

    在第一卦限内，$x = \sqrt{a^2-y^2}$．将其投影至 $yz$ 平面，区域由 $x \pm z = 0$ 及 $x>0$ 确定，即 $|z| \le \sqrt{a^2-y^2}$ 且 $0 < y < a$．
    \begin{equation}
        \d S = \sqrt{1 + x_y^2 + x_z^2} \d y \d z = \frac{a}{\sqrt{a^2-y^2}} \d y \d z
    \end{equation}
    \begin{equation}
        S = \int_0^a \d y \int_{-\sqrt{a^2-y^2}}^{\sqrt{a^2-y^2}} \frac{a}{\sqrt{a^2-y^2}} \d z = \int_0^a 2a \d y
    \end{equation}
    \begin{equation}
        \boxed{S = 2a^2}
    \end{equation}

    \ansitem{3}

    由对称性，总面积为第一卦限部分的 $8$ 倍．在第一卦限内，$x = \sqrt{a^2-y^2}$，投影至 $yz$ 平面的区域为 $y^2 + z^2 \le a^2$．
    \begin{equation}
        S = 8 \iint_{D_{yz}} \frac{a}{\sqrt{a^2-y^2}} \d y \d z = 8 \int_0^a \d y \int_0^{\sqrt{a^2-y^2}} \frac{a}{\sqrt{a^2-y^2}} \d z = 8 \int_0^a a \d y
    \end{equation}
    \begin{equation}
        \boxed{S = 8a^2}
    \end{equation}

    \ansitem{4}

    联立方程得交线为 $z^2 + 2az - 3a^2 = 0$，解得 $z=a$．全表面由球面部分 $S_1$ 和抛物面部分 $S_2$ 组成．
    \begin{equation}
        \begin{aligned}
            S_1 & = \iint_{x^2+y^2 \le 2a^2} \frac{\sqrt{3}a}{\sqrt{3a^2-x^2-y^2}} \d x \d y = 2\uppi \int_0^{\sqrt{2}a} \frac{\sqrt{3}ar}{\sqrt{3a^2-r^2}} \d r = 2\uppi a^2 (3-\sqrt{3})                                    \\
            S_2 & = \iint_{x^2+y^2 \le 2a^2} \sqrt{1 + \left( \frac{x}{a} \right)^2 + \left( \frac{y}{a} \right)^2} \d x \d y = \frac{2\uppi}{a} \int_0^{\sqrt{2}a} r \sqrt{a^2+r^2} \d r = \frac{2\uppi a^2}{3}(3\sqrt{3}-1)
        \end{aligned}
    \end{equation}
    \begin{equation}
        \boxed{S = S_1 + S_2 = \frac{16}{3} \uppi a^2}
    \end{equation}

    \ansitem{5}

    由 $x = y^2 + \frac{1}{2}z^2$ 得 $\d S = \sqrt{1 + 4y^2 + z^2} \d y \d z$．令 $y = \frac{1}{2}r \cos \theta, z = r \sin \theta$，雅可比行列式 $J = \frac{1}{2}r$．
    \begin{equation}
        S = \iint_{4y^2+z^2 \le 1} \sqrt{1+4y^2+z^2} \d y \d z = \int_0^{2\uppi} \d \theta \int_0^1 \frac{1}{2}r \sqrt{1+r^2} \d r
    \end{equation}
    \begin{equation}
        \boxed{S = \frac{\uppi}{3}(2\sqrt{2}-1)}
    \end{equation}

    \ansitem{6}

    锥面在 $Oxy$ 平面的投影 $D$ 由 $\sqrt{x^2+y^2} = \frac{x}{\sqrt{2}} + \sqrt{2}$ 确定，即椭圆 $\frac{(x-2)^2}{8} + \frac{y^2}{4} = 1$．
    \begin{equation}
        S = \iint_D \sqrt{1+z_x^2+z_y^2} \d x \d y = \iint_D \sqrt{2} \d x \d y = \sqrt{2} \cdot \uppi(2\sqrt{2})(2)
    \end{equation}
    \begin{equation}
        \boxed{S = 8\uppi}
    \end{equation}

    \ansitem{7}

    由参数方程 $\symbfit{r} = (r \cos \varphi, r \sin \varphi, h \varphi)$ 计算第一基本形式系数：
    \begin{equation}
        E = 1, \quad F = 0, \quad G = r^2 + h^2 \implies \d S = \sqrt{r^2 + h^2} \d r \d \varphi
    \end{equation}
    \begin{equation}
        S = \int_0^{2\uppi} \d \varphi \int_0^a \sqrt{r^2+h^2} \d r = 2\uppi \left[ \frac{r}{2}\sqrt{r^2+h^2} + \frac{h^2}{2} \ln (r + \sqrt{r^2+h^2}) \right]_0^a
    \end{equation}
    \begin{equation}
        \boxed{S = \uppi a \sqrt{a^2+h^2} + \uppi h^2 \ln \frac{a + \sqrt{a^2+h^2}}{h}}
    \end{equation}

    \ansitem{8}

    采用球坐标 $\rho = a \sin \varphi \sqrt{\sin 2\theta}$，其中 $\theta \in [0, \uppi/2] \cup [\uppi, 3\uppi/2]$．
    \begin{equation}
        \d S = \rho \sqrt{\rho^2 + \rho_\varphi^2 + \frac{\rho_\theta^2}{\sin^2 \varphi}} \sin \varphi \d \varphi \d \theta = a^2 \sin^2 \varphi \d \varphi \d \theta
    \end{equation}
    \begin{equation}
        S = a^2 \int_0^{\uppi} \sin^2 \varphi \d \varphi \left( \int_0^{\uppi/2} \d \theta + \int_{\uppi}^{3\uppi/2} \d \theta \right) = a^2 \cdot \frac{\uppi}{2} \cdot \uppi
    \end{equation}
    \begin{equation}
        \boxed{S = \frac{\uppi^2 a^2}{2}}
    \end{equation}
\end{solution}

\begin{problem}{曲面积分}{Surface-Integral-Calculations}
    计算下列曲面积分：
    \begin{enumerate}
        \item $\iint_{S}(x+y+z) \d S$，$S$：立方体 $0 \leqslant x \leqslant 1$、$0 \leqslant y \leqslant 1$、$0 \leqslant z \leqslant 1$ 的全表面；
        \item $\iint_{S} x y z \d S$，$S$：$x+y+z=1$ 在第一卦限部分；
        \item $\iint_{S}\left(x^{2}+y^{2}\right) \d S$，$S$：由 $z=\sqrt{x^{2}+y^{2}}$ 和 $z=1$ 所围成的立体表面；
        \item $\iint_{S}(x y+y z+z x) \d S$，$S$：锥面 $z=\sqrt{x^{2}+y^{2}}$ 被柱面 $x^{2}+y^{2}=2 a x$ ($a>0$) 所割下的那块曲面；
        \item $\iint_{S}(x^{4}-y^{4}+y^{2} z^{2}-x^{2} z^{2}+1) \d S$，$S$：圆锥 $z=\sqrt{x^{2}+y^{2}}$ 被柱面 $x^{2}+y^{2}=2 x$ 所截下的部分；
        \item $\iint_{S} \frac{\d S}{r^{2}}$，$S$：圆柱面 $x^{2}+y^{2}=R^{2}$ 界于平面 $z=0$ 及 $z=H$ 之间的部分，$r$ 是 $S$ 上的点到原点的距离；
        \item $\iint_{S}|x y z| \d S$，$S$ 为曲面 $z=x^{2}+y^{2}$ 介于二平面 $z=0$ 和 $z=1$ 间的部分．
    \end{enumerate}
\end{problem}

\begin{solution}
    \ansitem{1}

    根据对称性，可知：
    \begin{equation}
        \iint_{S} x \d S = \iint_{S} y \d S = \iint_{S} z \d S
    \end{equation}
    考察 $\iint_{S} z \d S$，其由六个面组成．在 $z=1$ 面上，积分值为 $1$；在 $z=0$ 面上，积分值为 $0$；在其余四个侧面上，$z$ 的积分为 $\int_{0}^{1} z \d z = \frac{1}{2}$．故：
    \begin{equation}
        \iint_{S} z \d S = 1 + 0 + 4 \times \frac{1}{2} = 3
    \end{equation}
    因此，原积分值为：
    \begin{equation}
        \boxed{\iint_{S}(x+y+z) \d S = 9}
    \end{equation}

    \ansitem{2}

    曲面 $S$ 在 $xOy$ 面上的投影区域为 $D = \{(x, y) \mid 0 \leqslant x \leqslant 1, 0 \leqslant y \leqslant 1-x\}$．由于 $z = 1-x-y$，故 $\d S = \sqrt{1 + z_x^2 + z_y^2} \d x \d y = \sqrt{3} \d x \d y$．
    \begin{equation}
        \begin{aligned}
            \iint_{S} x y z \d S & = \sqrt{3} \iint_{D} x y(1-x-y) \d x \d y                                \\
                                 & = \sqrt{3} \int_{0}^{1} x \d x \int_{0}^{1-x} (y - xy - y^2) \d y        \\
                                 & = \sqrt{3} \int_{0}^{1} \frac{1}{6} x(1-x)^3 \d x = \frac{\sqrt{3}}{120}
        \end{aligned}
    \end{equation}
    \begin{equation}
        \boxed{\iint_{S} x y z \d S = \frac{\sqrt{3}}{120}}
    \end{equation}

    \ansitem{3}

    表面 $S$ 由锥面 $S_1: z = \sqrt{x^2+y^2}$ ($x^2+y^2 \leqslant 1$) 和顶面 $S_2: z=1$ ($x^2+y^2 \leqslant 1$) 组成．
    对于 $S_1$，$\d S = \sqrt{2} \d x \d y$：
    \begin{equation}
        \iint_{S_1} (x^2+y^2) \d S = \sqrt{2} \iint_{x^2+y^2 \leqslant 1} (x^2+y^2) \d x \d y = \sqrt{2} \int_{0}^{2\uppi} \d \theta \int_{0}^{1} r^3 \d r = \frac{\uppi \sqrt{2}}{2}
    \end{equation}
    对于 $S_2$，$\d S = \d x \d y$：
    \begin{equation}
        \iint_{S_2} (x^2+y^2) \d S = \iint_{x^2+y^2 \leqslant 1} (x^2+y^2) \d x \d y = \frac{\uppi}{2}
    \end{equation}
    \begin{equation}
        \boxed{\iint_{S} (x^2+y^2) \d S = \frac{\uppi}{2}(1+\sqrt{2})}
    \end{equation}

    \ansitem{4}

    在锥面 $z=\sqrt{x^2+y^2}$ 上，$\d S = \sqrt{2} \d x \d y$．投影区域为 $D: x^2+y^2 \leqslant 2ax$．利用对称性 $\iint_D xy \d A = 0$：
    \begin{equation}
        \begin{aligned}
            \iint_{S} (xy+yz+zx) \d S & = \sqrt{2} \iint_{D} (xy + (x+y)\sqrt{x^2+y^2}) \d x \d y                                                                \\
                                      & = \sqrt{2} \iint_{D} x\sqrt{x^2+y^2} \d x \d y                                                                           \\
                                      & = \sqrt{2} \int_{-\uppi/2}^{\uppi/2} \cos \theta \d \theta \int_{0}^{2a\cos \theta} r^3 \d r = \frac{64\sqrt{2}}{15} a^4
        \end{aligned}
    \end{equation}
    \begin{equation}
        \boxed{\iint_{S} (xy+yz+zx) \d S = \frac{64\sqrt{2}}{15} a^4}
    \end{equation}

    \ansitem{5}

    注意到被积函数可化为 $(x^2-y^2)(x^2+y^2-z^2) + 1$．在锥面 $z^2 = x^2+y^2$ 上，前项为 $0$．故：
    \begin{equation}
        \iint_{S} (x^4-y^4+y^2z^2-x^2z^2+1) \d S = \iint_{S} 1 \d S = \sqrt{2} \iint_{x^2+y^2 \leqslant 2x} \d x \d y = \sqrt{2} \uppi
    \end{equation}
    \begin{equation}
        \boxed{\iint_{S} (x^4-y^4+y^2z^2-x^2z^2+1) \d S = \sqrt{2}\uppi}
    \end{equation}

    \ansitem{6}

    在柱面上 $x^2+y^2=R^2$，则 $r^2 = R^2+z^2$．采用柱坐标，$x=R\cos\theta, y=R\sin\theta, z=z$，面积元 $\d S = R \d \theta \d z$．
    \begin{equation}
        \iint_{S} \frac{\d S}{r^2} = \int_{0}^{H} \d z \int_{0}^{2\uppi} \frac{R}{R^2+z^2} \d \theta = 2\uppi R \int_{0}^{H} \frac{1}{R^2+z^2} \d z = 2\uppi \arctan \frac{H}{R}
    \end{equation}
    \begin{equation}
        \boxed{\iint_{S} \frac{\d S}{r^2} = 2\uppi \arctan \frac{H}{R}}
    \end{equation}

    \ansitem{7}

    由于曲面方程为 $z=x^{2}+y^{2}$，其在 $xOy$ 面上的投影区域为 $D: x^{2}+y^{2} \leqslant 1$．面积元为：
    \begin{equation}
        \d S = \sqrt{1+\left(\frac{\partial z}{\partial x}\right)^{2}+\left(\frac{\partial z}{\partial y}\right)^{2}} \d x \d y = \sqrt{1+4\left(x^{2}+y^{2}\right)} \d x \d y
    \end{equation}
    将积分转换至极坐标系下，令 $x=r \cos \theta$，$y=r \sin \theta$：
    \begin{equation}
        \begin{aligned}
            \iint_{S}|x y z| \d S & = \iint_{D}\left|x y\left(x^{2}+y^{2}\right)\right| \sqrt{1+4\left(x^{2}+y^{2}\right)} \d x \d y                           \\
                                  & = \int_{0}^{2 \uppi}|\sin \theta \cos \theta| \d \theta \int_{0}^{1} r^{2} \cdot r^{2} \cdot \sqrt{1+4 r^{2}} \cdot r \d r \\
                                  & = \frac{1}{2} \int_{0}^{2 \uppi}|\sin 2 \theta| \d \theta \int_{0}^{1} r^{5} \sqrt{1+4 r^{2}} \d r
        \end{aligned}
    \end{equation}
    计算角度部分积分：
    \begin{equation}
        \int_{0}^{2 \uppi}|\sin 2 \theta| \d \theta = 4 \int_{0}^{\uppi / 2} \sin 2 \theta \d \theta = 4 \left[ -\frac{1}{2} \cos 2 \theta \right]_{0}^{\uppi / 2} = 4
    \end{equation}
    对于径向部分，令 $u=\sqrt{1+4 r^{2}}$，则 $r^{2}=\frac{u^{2}-1}{4}$，$r \d r = \frac{1}{4} u \d u$：
    \begin{equation}
        \begin{aligned}
            \int_{0}^{1} r^{5} \sqrt{1+4 r^{2}} \d r & = \int_{1}^{\sqrt{5}}\left(\frac{u^{2}-1}{4}\right)^{2} \cdot u \cdot \frac{1}{4} u \d u        \\
                                                     & = \frac{1}{64} \int_{1}^{\sqrt{5}}\left(u^{6}-2 u^{4}+u^{2}\right) \d u                         \\
                                                     & = \frac{1}{64}\left[\frac{1}{7} u^{7}-\frac{2}{5} u^{5}+\frac{1}{3} u^{3}\right]_{1}^{\sqrt{5}} \\
                                                     & = \frac{125 \sqrt{5}-1}{840}
        \end{aligned}
    \end{equation}
    综合上述计算结果，得：
    \begin{equation}
        \boxed{\iint_{S}|x y z| \d S = \frac{125 \sqrt{5}-1}{420}}
    \end{equation}
\end{solution}

\begin{problem}{利用对称性计算曲面积分}{Surface-Integral-Symmetry}
    利用对称性计算下列曲面积分：
    \begin{enumerate}
        \item $\iint_{S} (x^2 + y^2) \d S$，其中 $S$ 为球面 $x^2 + y^2 + z^2 = R^2$；
        \item $\iint_{S} (x + y + z) \d S$，其中 $S$ 为半球面 $x^2 + y^2 + z^2 = a^2$ ($z \geqslant 0$)．
    \end{enumerate}
\end{problem}

\begin{solution}
    \ansitem{1}

    由于曲面 $S$ 关于坐标面具有轮换对称性，故有
    \begin{equation}
        \iint_{S} x^2 \d S = \iint_{S} y^2 \d S = \iint_{S} z^2 \d S. \label{eq:symmetry-1}
    \end{equation}
    在球面 $S$ 上满足 $x^2 + y^2 + z^2 = R^2$，则
    \begin{equation}
        \iint_{S} (x^2 + y^2 + z^2) \d S = \iint_{S} R^2 \d S = R^2 \cdot 4\uppi R^2 = 4\uppi R^4. \label{eq:total-integral}
    \end{equation}
    结合 \zcref{eq:symmetry-1} 与 \zcref{eq:total-integral} 可得
    \begin{equation}
        3 \iint_{S} x^2 \d S = 4\uppi R^4 \implies \iint_{S} x^2 \d S = \frac{4}{3}\uppi R^4.
    \end{equation}
    因此，所求积分为
    \begin{equation}
        \iint_{S} (x^2 + y^2) \d S = 2 \iint_{S} x^2 \d S = \frac{8}{3}\uppi R^4.
    \end{equation}
    \begin{equation}
        \boxed{\iint_{S} (x^2 + y^2) \d S = \frac{8}{3}\uppi R^4}
    \end{equation}

    \ansitem{2}

    由于曲面 $S$ 关于 $yz$ 平面及 $xz$ 平面对称，且被积函数 $x$ 与 $y$ 分别为关于 $x$ 与 $y$ 的奇函数，故
    \begin{equation}
        \iint_{S} x \d S = 0, \quad \iint_{S} y \d S = 0.
    \end{equation}
    从而原积分简化为
    \begin{equation}
        \iint_{S} (x + y + z) \d S = \iint_{S} z \d S.
    \end{equation}
    采用球面坐标系，令 $z = a \cos \theta$，面积元素 $\d S = a^2 \sin \theta \d \theta \d \varphi$．对于上半球面 $S$，积分区域为 $0 \leqslant \varphi \leqslant 2\uppi$，$0 \leqslant \theta \leqslant \frac{\uppi}{2}$，则
    \begin{equation}
        \begin{aligned}
            \iint_{S} z \d S & = \int_{0}^{2\uppi} \d \varphi \int_{0}^{\frac{\uppi}{2}} (a \cos \theta) (a^2 \sin \theta) \d \theta \\
                             & = 2\uppi a^3 \int_{0}^{\frac{\uppi}{2}} \sin \theta \cos \theta \d \theta                             \\
                             & = 2\uppi a^3 \left[ \frac{1}{2} \sin^2 \theta \right]_{0}^{\frac{\uppi}{2}} = \uppi a^3.
        \end{aligned}
    \end{equation}
    \begin{equation}
        \boxed{\iint_{S} (x + y + z) \d S = \uppi a^3}
    \end{equation}
\end{solution}

\begin{problem}{面积投影关系}{Surface-Area-Projection}
    设 $G$ 是平面 $Ax + By + Cz + D = 0$ ($C \neq 0$) 上的一个有界闭区域，它在 $Oxy$ 平面上的投影是 $G_1$，试证：$\frac{G \text{ 的面积}}{G_1 \text{ 的面积}} = \sqrt{\frac{A^2 + B^2 + C^2}{C^2}}$．
\end{problem}

\begin{myproof}

    将平面方程表示为
    \begin{equation}
        z = -\frac{A}{C}x - \frac{B}{C}y - \frac{D}{C}.
    \end{equation}
    计算其偏导数
    \begin{equation}
        \frac{\partial z}{\partial x} = -\frac{A}{C}, \quad \frac{\partial z}{\partial y} = -\frac{B}{C}.
    \end{equation}
    区域 $G$ 的面积 $S_G$ 可通过曲面积分计算
    \begin{equation}
        \begin{aligned}
            S_G & = \iint_{G_1} \sqrt{1 + \left( \frac{\partial z}{\partial x} \right)^2 + \left( \frac{\partial z}{\partial y} \right)^2} \d x \d y \\
                & = \iint_{G_1} \sqrt{1 + \frac{A^2}{C^2} + \frac{B^2}{C^2}} \d x \d y                                                               \\
                & = \frac{\sqrt{A^2 + B^2 + C^2}}{|C|} \iint_{G_1} \d x \d y.
        \end{aligned}
    \end{equation}
    记 $G_1$ 的面积为 $S_{G_1}$，则有
    \begin{equation}
        \frac{S_G}{S_{G_1}} = \frac{\sqrt{A^2 + B^2 + C^2}}{|C|} = \sqrt{\frac{A^2 + B^2 + C^2}{C^2}}.
    \end{equation}
\end{myproof}

\begin{problem}{曲面积分求质量}{Surface-Integral-Mass}
    求抛物面壳 $z = \frac{1}{2}(x^2 + y^2)$ ($0 \le z \le 1$) 的质量，其各点的密度为 $\rho = z$．
\end{problem}

\begin{solution}
    曲面 $\Sigma$ 在 $Oxy$ 平面上的投影区域为 $D_{xy} = \{ (x, y) \mid x^2 + y^2 \le 2 \}$．计算偏导数
    \begin{equation}
        \frac{\partial z}{\partial x} = x, \quad \frac{\partial z}{\partial y} = y.
    \end{equation}
    面积元素为
    \begin{equation}
        \d S = \sqrt{1 + x^2 + y^2} \d x \d y.
    \end{equation}
    质量 $M$ 的积分为
    \begin{equation}
        \begin{aligned}
            M & = \iint_{\Sigma} \rho \d S = \iint_{D_{xy}} \frac{1}{2}(x^2 + y^2) \sqrt{1 + x^2 + y^2} \d x \d y \\
              & = \int_0^{2\uppi} \d \theta \int_0^{\sqrt{2}} \frac{1}{2} r^2 \sqrt{1 + r^2} \cdot r \d r         \\
              & = \uppi \int_0^{\sqrt{2}} r^3 \sqrt{1 + r^2} \d r.
        \end{aligned}
    \end{equation}
    令 $t = \sqrt{1 + r^2}$，则 $r^2 = t^2 - 1$，$r \d r = t \d t$．代入得
    \begin{equation}
        \begin{aligned}
            M & = \uppi \int_1^{\sqrt{3}} (t^2 - 1) t^2 \d t                                       \\
              & = \uppi \left[ \frac{1}{5} t^5 - \frac{1}{3} t^3 \right]_1^{\sqrt{3}}              \\
              & = \uppi \left( \frac{9\sqrt{3}}{5} - \sqrt{3} - \frac{1}{5} + \frac{1}{3} \right).
        \end{aligned}
    \end{equation}
    整理得
    \begin{equation}
        \boxed{M = \frac{2\uppi}{15} (6\sqrt{3} + 1)}
    \end{equation}
\end{solution}

\begin{problem}{转动惯量}{Moment-Of-Inertia-Spherical-Shell}
    一个半径为 $R$ 的均匀球壳（密度为 $\rho$）绕其直径旋转，求它的转动惯量．
\end{problem}

\begin{solution}
    建立空间直角坐标系，使球心位于原点，旋转轴为 $z$ 轴．球壳方程为 $x^2 + y^2 + z^2 = R^2$．
    根据转动惯量定义，有
    \begin{equation}
        \label{eq:Inertia-Def-1}
        I = \iint_{S} (x^2 + y^2) \rho \d S.
    \end{equation}
    采用球面坐标系，令 $x = R \sin \theta \cos \varphi$，$y = R \sin \theta \sin \varphi$，$z = R \cos \theta$．面积元为 $\d S = R^2 \sin \theta \d \theta \d \varphi$．
    将坐标代入 \zcref{eq:Inertia-Def-1} 可得
    \begin{equation}
        \begin{aligned}
            I & = \int_{0}^{2\uppi} \d \varphi \int_{0}^{\uppi} \rho (R^2 \sin^2 \theta) (R^2 \sin \theta) \d \theta \\
              & = 2\uppi \rho R^4 \int_{0}^{\uppi} \sin^3 \theta \d \theta                                           \\
              & = 2\uppi \rho R^4 \int_{0}^{\uppi} (1 - \cos^2 \theta) \sin \theta \d \theta.
        \end{aligned}
    \end{equation}
    计算积分项
    \begin{equation}
        \int_{0}^{\uppi} (1 - \cos^2 \theta) \sin \theta \d \theta = \left[ -\cos \theta + \frac{1}{3} \cos^3 \theta \right]_{0}^{\uppi} = \frac{4}{3}.
    \end{equation}
    故最终结果为
    \begin{equation}
        I = \boxed{\frac{8}{3} \uppi \rho R^4}
    \end{equation}
\end{solution}

\begin{problem}{引力}{Gravitational-Force-Conical-Surface}
    一个密度为 $\rho$ 的均匀截锥面 $z = \sqrt{x^2 + y^2}$ ($0 < a \leqslant z \leqslant b$)，求它对于处在锥顶的质量为 $m$ 的质点的引力．
\end{problem}

\begin{solution}
    设锥顶位于原点 $(0,0,0)$．由对称性可知，引力 $\symbfit{F}$ 在 $x$ 与 $y$ 方向的分量为零，仅存在 $z$ 方向分量 $F_z$．
    设锥面上任一点为 $(x, y, z)$，其到原点的距离为
    \begin{equation}
        r = \sqrt{x^2 + y^2 + z^2} = \sqrt{z^2 + z^2} = \sqrt{2} z.
    \end{equation}
    该点处位置矢量与 $z$ 轴夹角的余弦为 $\cos \gamma = \frac{z}{r} = \frac{1}{\sqrt{2}}$．
    面积元为 $\d S = \sqrt{1 + z_x^2 + z_y^2} \d x \d y = \sqrt{2} \d x \d y$．在 $xy$ 平面上采用极坐标，有 $\d S = \sqrt{2} r_{xy} \d r_{xy} \d \theta$．由于在锥面上 $z = r_{xy}$，故 $\d S = \sqrt{2} z \d z \d \theta$．
    根据万有引力定律，引力分量 $F_z$ 为
    \begin{equation}
        \begin{aligned}
            F_z & = \iint_{S} \frac{G m \rho \d S}{r^2} \cos \gamma                                                               \\
                & = \int_{0}^{2\uppi} \d \theta \int_{a}^{b} \frac{G m \rho (\sqrt{2} z \d z)}{(\sqrt{2} z)^2} \frac{1}{\sqrt{2}} \\
                & = \int_{0}^{2\uppi} \d \theta \int_{a}^{b} \frac{G m \rho}{2 z} \d z.
        \end{aligned}
    \end{equation}
    计算上述积分得
    \begin{equation}
        F_z = 2\uppi \cdot \frac{G m \rho}{2} \ln \frac{b}{a} = \uppi G m \rho \ln \frac{b}{a}.
    \end{equation}
    因此，引力矢量为
    \begin{equation}
        \symbfit{F} = \boxed{\uppi G m \rho \ln \frac{b}{a} \symbfit{k}}
    \end{equation}
\end{solution}

\endinput
